\documentclass[11pt]{article}
\usepackage[a4paper,margin=1in]{geometry}
\usepackage{amsmath,amssymb,amsthm,mathtools}
\usepackage{graphicx}
\usepackage[hidelinks]{hyperref}
\usepackage{microtype}

\newtheorem{theorem}{Theorem}
\newtheorem{lemma}[theorem]{Lemma}
\newtheorem{proposition}[theorem]{Proposition}
\newcommand{\T}{\mathbb T}
\newcommand{\D}{\mathbb D}
\newcommand{\C}{\mathbb C}
\newcommand{\R}{\mathbb R}
\newcommand{\G}{\mathcal G}

\title{The weighted two-atom Cauchy dual classification}
\author{Candidate solution packet for arXiv:2510.14004}
\date{13 August 2026}

\begin{document}
\maketitle

\begin{center}
\fbox{\parbox{0.9\textwidth}{\textbf{Status.} Candidate full solution, likely valid, awaiting specialist review.  The new direction is non-subnormality for arbitrary positive weights at two distinct non-antipodal points.  The antipodal direction is due to Chavan--Ghara--Reza.}}
\end{center}

\section{Source question and answer}

Khasnis, Phatak, and Sholapurkar prove the classification for the unweighted
measure $\delta_{\zeta_1}+\delta_{\zeta_2}$ and close their paper with the
following weighted extension \cite{KPS} (source PDF, p.~25, \S4):

\begin{center}
\includegraphics[width=0.98\textwidth]{figures/open_problem_crop.png}
\end{center}

In their words, they expect the main theorem for
$\mu=c_1\delta_{\zeta_1}+c_2\delta_{\zeta_2}$ with $c_1,c_2>0$.
The following gives precisely that extension.

\begin{theorem}[Weighted two-atom classification]\label{thm:main}
Let $c_1,c_2>0$ and let $\zeta_1,\zeta_2\in\T$ be distinct.  On the
Dirichlet-type space
\[
   D(c_1\delta_{\zeta_1}+c_2\delta_{\zeta_2}),
\]
the Cauchy dual of multiplication by $z$ is subnormal if and only if
$\zeta_2=-\zeta_1$.
\end{theorem}

The ``if'' direction, including arbitrary nonnegative weights, is
Theorem~2.3 of Chavan--Ghara--Reza \cite{CGR}.  The proof below establishes
the converse for all positive weights.

\section{Proof intuition}

The Dirichlet space is first written as a rank-two de Branges--Rovnyak
space $H(B)$ whose denominator is the outer Fej\'er--Riesz factor
$q=(z-\alpha_1)(z-\alpha_2)$.  Chavan--Ghara--Reza turn subnormality into
positivity conditions involving the numerator Gram function
\[
 \G(z,w)=\sum_{j=1}^2p_j(z)\overline{p_j(w)}.
\]
For unequal weights, the two outer roots cannot lie on one ray.  The
simple-pole criterion would therefore force $\G(\alpha_1,\alpha_2)=0$.
After rotating by half the angle between the atoms, this vanishing has a
single realness consequence.  Eliminating the remaining root angle from
that consequence and the root equation gives a completely factored,
nonzero resultant.

Equal weights have an additional reflection symmetry.  The outer roots
are either a nonreal conjugate pair, two distinct positive roots, or a
double positive root.  In the first case the cross Gram value is directly
nonzero.  In the second, a $2\times2$ positivity minor becomes negative
for sufficiently high difference order.  In the double-root case the same
minor is affine in that order and has a strictly negative leading
coefficient.  These three arguments cover every non-antipodal
configuration.

\section{The rank-two model and criterion}

Rotation and interchange of the atoms preserve the question.  For the
non-antipodal direction we may therefore write
\[
  \mu=c_1\delta_1+c_2\delta_\xi,\qquad
  \xi=e^{2i\gamma},\qquad 0<\gamma<\frac\pi2.
\]
Set $C=c_1+c_2$.  On $\T$, let
\begin{align*}
 F(z)&=|z-1|^2|z-\xi|^2+c_1|z-\xi|^2+c_2|z-1|^2.
\end{align*}
The rank-two construction of \cite{CGR}, also recalled in \cite{KPS},
gives a monic outer factor
\[
 q(z)=z^2-(u+v\xi)z+b\xi=(z-\alpha_1)(z-\alpha_2),
 \qquad |\alpha_j|>1,\quad b>1,
\]
such that $F=b^{-1}|q|^2$ on $\T$.  Comparing coefficients gives
\begin{equation}\label{eq:uv}
 u=\frac{b\{b(c_1+2)-(c_2+2)\}}{b^2-1},\qquad
 v=\frac{b\{b(c_2+2)-(c_1+2)\}}{b^2-1},
\end{equation}
and
\begin{equation}\label{eq:physical}
 b^2+1+u^2+v^2+2uv\cos(2\gamma)
 =2b\{C+2+\cos(2\gamma)\}.
\end{equation}

There is a row rational Schur function
$B=(p_1/q,p_2/q)$ for which $D(\mu)=H(B)$ isometrically.  We use the
coefficient convention
\[
 \G(z,w)=\sum_{j=1}^2p_j(z)\overline{p_j(w)}
        =(z,z^2)A(\overline w,\overline w^2)^T.
\]
The kernel calculation in the rank-two construction simplifies to
\begin{equation}\label{eq:gram-general}
\G(z,w)=z\overline w\bigl[A_0-b(c_2+c_1\overline\xi)z
 -b(c_2+c_1\xi)\overline w+A_2z\overline w\bigr],
\end{equation}
where $A_0,A_2\in\R$ and
\begin{equation}\label{eq:diag-sum}
 A_0+A_2=2bC.
\end{equation}
One quick check on \eqref{eq:gram-general} is the boundary identity
\begin{equation}\label{eq:boundary}
 \G(z,z)=b\{c_1|z-\xi|^2+c_2|z-1|^2\},\qquad z\in\T.
\end{equation}
This also fixes the orientation of the two off-diagonal coefficients.

We use two consequences of \cite{CGR}.  If $\alpha_1,\alpha_2$ are
distinct and $\alpha_1\overline{\alpha_2}\notin[1,\infty)$, their
Corollary~4.3 says that subnormality is equivalent to
\begin{equation}\label{eq:crosszero}
 \G(\alpha_1,\alpha_2)=0.
\end{equation}
More generally, their Theorem~2.1 requires, for every integer $\ell\ge1$,
formal positivity of
\begin{equation}\label{eq:cgrmatrix}
 \sum_{r,t=1}^2
 \frac{\G(\alpha_r,\alpha_t)}{a_r\overline{a_t}}
 \left(1-\frac1{\alpha_r\overline{\alpha_t}}\right)^\ell
 \left(\!\left(\frac1{\alpha_r^{m+2}
 \overline{\alpha_t}^{,n+2}}\right)\!\right)_{m,n\ge0},
 \quad a_1=\alpha_1-\alpha_2=-a_2.
\end{equation}

\section{Unequal weights}

Assume $c_1\ne c_2$ and put
\[
 d=c_1-c_2,\quad H=\cos^2\gamma,\quad
 m=u+v=\frac{b(C+4)}{b+1},\quad
 n=v-u=-\frac{bd}{b-1}.
\]
Equation \eqref{eq:physical} becomes
\begin{equation}\label{eq:physical-mn}
 b^2+1+m^2H+n^2(1-H)=2b(C+1+2H).
\end{equation}

Write the outer roots in the form
\begin{equation}\label{eq:root-param}
 \alpha_1=e^{i\gamma}R e^{i\phi},\qquad
 \alpha_2=e^{i\gamma}\frac bR e^{-i\phi},\qquad R>0.
\end{equation}
Their sum gives
\begin{equation}\label{eq:PQ}
 \left(R+\frac bR\right)\cos\phi=m\sqrt H,\qquad
 \left(R-\frac bR\right)\sin\phi=n\sqrt{1-H}.
\end{equation}
The first right side is positive and the second is nonzero.  Hence
$\cos\phi\ne0$ and $\sin\phi\ne0$, so
$\alpha_1\overline{\alpha_2}=be^{2i\phi}$ is nonreal.  In particular the
roots are distinct and the simple-pole criterion applies.  Put
$\chi=\sin^2\phi$.  Eliminating $R$ from \eqref{eq:PQ} yields
\begin{equation}\label{eq:Rchi}
 \mathcal R(\chi):=
 m^2H\chi-n^2(1-H)(1-\chi)-4b\chi(1-\chi)=0.
\end{equation}

We now rule out \eqref{eq:crosszero}.  If it held, then by
\eqref{eq:gram-general}--\eqref{eq:diag-sum},
\begin{equation}\label{eq:ratio-real}
 \frac{A_2}{b}=
 \frac{W-2C}{\alpha_1\overline{\alpha_2}-1}\in\R,
 \quad
 W=(c_2+c_1\overline\xi)\alpha_1
 +(c_2+c_1\xi)\overline{\alpha_2}.
\end{equation}
Using \eqref{eq:PQ}, the imaginary part in \eqref{eq:ratio-real}
reduces exactly to
\begin{equation}\label{eq:Echi}
 \mathcal E(\chi):=
 d^2(1-H)(1-\chi)-C(C+4)H\chi+4C\chi(1-\chi)=0.
\end{equation}
For completeness, one intermediate form is
\[
 W=e^{i\phi}\left\{
 \frac{bC(C+4)H}{(b+1)\cos\phi}
 +i\frac{bd^2(1-H)}{(b-1)\sin\phi}\right\},
\]
from which \eqref{eq:Echi} follows by one cross multiplication.

It remains to show that \eqref{eq:Rchi} and \eqref{eq:Echi} cannot share
a root.  Let $L=b^2-(C+2)b+1$.  First, $uv\ne b$.  Indeed, otherwise
\eqref{eq:physical-mn} and $m^2-n^2=4b$ imply
\[
 0=\frac{L^2}{(b+1)^2}.
\]
Then $C=(b-1)^2/b$, $m=b+1$, and $m^2-n^2=4b$ gives
$|d|=C$, impossible because $c_1,c_2>0$.  Hence
\eqref{eq:physical-mn} determines
\begin{equation}\label{eq:Hsolve}
 H=\frac{2b(C+1)-b^2-1-n^2}{4(uv-b)}.
\end{equation}
Also $L\ne0$: if $L=0$, substituting $C=(b-1)^2/b$ into
\eqref{eq:physical-mn} gives
\[
 \bigl((b-1)^2-n^2\bigr)(H-1)=0.
\]
Since $0<H<1$, this again forces $|d|=C$.

A direct quadratic resultant, with \eqref{eq:Hsolve} substituted, is
\begin{equation}\label{eq:resultant}
 \operatorname{Res}_{\chi}(\mathcal R,\mathcal E)
 =\frac{H b^3d^2(d^2-C^2)(C+4)L^5}
 {(uv-b)^2(b-1)^2(b+1)^6}.
\end{equation}
Every factor on the right is nonzero: in particular
$d^2-C^2=-4c_1c_2<0$.  Thus \eqref{eq:Rchi} and
\eqref{eq:Echi} have no common root, contradicting
\eqref{eq:ratio-real}.  Therefore $\G(\alpha_1,\alpha_2)\ne0$, and
\eqref{eq:crosszero} proves non-subnormality for unequal weights.

\section{Equal weights: simple roots}

Now let $c_1=c_2=c>0$.  Rotate the support symmetrically to
$\{e^{-i\gamma},e^{i\gamma}\}$ and write $h=\cos\gamma\in(0,1)$.  The
outer factor and the physical coefficient relation are
\begin{equation}\label{eq:eq-q}
 q(z)=z^2-2ahz+b,\qquad a=\frac{b(c+2)}{b+1},
\end{equation}
\begin{equation}\label{eq:eq-phys}
 b^2+1+4a^2h^2=2b(2c+1+2h^2).
\end{equation}
The possibility $a^2=b$ is excluded: together with
\eqref{eq:eq-phys} it would give both
$b+b^{-1}=4c+2$ and $b+b^{-1}=c^2+4c+2$, hence $c=0$.
Thus
\begin{equation}\label{eq:h2}
 h^2=\frac{(b+1)^2(b^2-4bc-2b+1)}
 {4b(b^2-bc^2-4bc-2b+1)}.
\end{equation}

The numerator Gram function is now
\begin{equation}\label{eq:equal-gram}
 \G(z,w)=z\overline w\{A_{11}+A_{12}(z+\overline w)
 +A_{22}z\overline w\},
\end{equation}
where
\begin{equation}\label{eq:Aequal}
 A_{12}=-2bch,\qquad A_{11}+A_{22}=4bc,\qquad
 A_{22}=-\frac{b^3-4b^2c-3b^2+3b-1}{b+1}.
\end{equation}
These identities follow by the same rank-two kernel calculation as
\eqref{eq:gram-general}.  An exact derivation is included in
\texttt{code/equal\_pair\_algebra.py}.

Suppose first that $q$ has nonreal roots.  They are
$r=\sqrt b e^{i\phi}$ and $\overline r$, with
$\cos\phi=ah/\sqrt b$ and $\sin\phi\ne0$.  Up to the nonzero factor
$r\overline{\overline r}$, the cross Gram value is
\begin{equation}\label{eq:Econj}
 E=A_{11}-4bch\,r+A_{22}r^2.
\end{equation}
If $E=0$, its imaginary part gives $aA_{22}=2bc$.  Under that relation,
its real part is
\[
 E=\frac{2bc}{a}(2a-1-b).
\]
Consequently $E=0$ would imply
\begin{equation}\label{eq:Leq}
 b^2-2b(c+1)+1=0.
\end{equation}
Substitution in \eqref{eq:eq-phys}, or in \eqref{eq:h2}, gives $h^2=1$,
a contradiction.  Thus the cross Gram value is nonzero.  Since
$r\overline{\overline r}=r^2\notin[1,\infty)$, the simple-pole criterion
gives non-subnormality.

Suppose next that $q$ has two distinct real roots $1<r_1<r_2$.  Then
$r_1r_2=b$ and $r_1+r_2=2ah$.  The cross Gram bracket factors as
\begin{equation}\label{eq:Ecommon}
 \frac{\G(r_1,r_2)}{r_1r_2}
 =-\frac{(b^2-2bc-2b+1)^3}
 {(b+1)(b^2-bc^2-4bc-2b+1)}.
\end{equation}
The denominator is nonzero because
\[
 b^2-bc^2-4bc-2b+1=-\frac{(b+1)^2}{b}(a^2-b),
\]
and distinct real roots imply $a^2h^2>b$, hence $a^2>b$.  The numerator
cannot vanish, since its vanishing is \eqref{eq:Leq}, which again forces
$h^2=1$.  Therefore \eqref{eq:Ecommon} is nonzero.

For these positive roots the leading $2\times2$ minor of
\eqref{eq:cgrmatrix} has the sign of
\begin{multline}\label{eq:common-minor}
 \G(r_1,r_1)\G(r_2,r_2)
 \{(1-r_1^{-2})(1-r_2^{-2})\}^{\ell}\\
 -|\G(r_1,r_2)|^2\{1-(r_1r_2)^{-1}\}^{2\ell}.
\end{multline}
But
\[
 Q:=\frac{(r_1^2-1)(r_2^2-1)}{(r_1r_2-1)^2}<1,
 \qquad
 (r_1r_2-1)^2-(r_1^2-1)(r_2^2-1)=(r_2-r_1)^2.
\]
The ratio of the positive term in \eqref{eq:common-minor} to the absolute
value of the negative term is a fixed finite constant times $Q^\ell$.
It is therefore less than one for all sufficiently large $\ell$, giving a
negative minor and non-subnormality.

\section{Equal weights: the double root}

It remains to consider
\[
 q(z)=(z-R)^2,\qquad R>1.
\]
Here $b=R^2$, $h=R/a$, and \eqref{eq:eq-phys} is equivalent to
\begin{equation}\label{eq:repeat-relation}
 \left(R+R^{-1}\right)^2=\frac{4(c+2)^2}{c+4}.
\end{equation}
Write the numerator row as
\[
 p(z)=uz+vz^2,\qquad
 \langle u,u\rangle=A_{11},\quad
 \langle u,v\rangle=A_{12},\quad
 \langle v,v\rangle=A_{22}.
\]
The Taylor coefficient row of $B=p/q$ is, for $n\ge1$,
\begin{equation}\label{eq:Bn}
 B_n=R^{-(n+1)}\{nu+(n-1)Rv\}.
\end{equation}
The general coefficient criterion in Theorem~3.5 of \cite{CGR} requires
formal positivity, for every $\ell\ge1$, of
\begin{equation}\label{eq:repeat-matrix}
 D_\ell=\left(\!\left(
 \sum_{s=0}^{\ell}(-1)^s\binom{\ell}{s}
 B_{m+1+s}B_{n+1+s}^*
 \right)\!\right)_{m,n\ge0}.
\end{equation}

Insert \eqref{eq:Bn} and take the leading $2\times2$ minor.  After
removing positive scalar factors and a positive diagonal similarity, its
determinant is affine in $\ell$.  Its coefficient of $\ell$ is
\begin{equation}\label{eq:repeat-slope}
 -\frac{(R^2-1)^2}{R^6}
 \bigl(A_{11}+2RA_{12}+R^2A_{22}\bigr)^2.
\end{equation}
The factor being squared is not zero.  Indeed, by \eqref{eq:Aequal},
with $b=R^2$,
\begin{align*}
 A_{22}&=-\frac{R^6-4R^4c-3R^4+3R^2-1}{R^2+1},\\
 A_{11}&=4R^2c-A_{22},\qquad
 A_{12}=-\frac{2Rc(R^2+1)}{c+2},
\end{align*}
and
\begin{equation}\label{eq:J-K}
 A_{11}+2RA_{12}+R^2A_{22}
 =\frac{K}{(R^2+1)(c+2)},
\end{equation}
where
\[
 K=-(c+2)(R^8+1)+4(c^2+2c+2)(R^6+R^2)-(14c+12)R^4.
\]
Put $Y=R+R^{-1}$.  Using \eqref{eq:repeat-relation},
\begin{align*}
 \frac K{R^4}
 &=-(c+2)(Y^4-4Y^2+2)
   +4(c^2+2c+2)(Y^2-2)-(14c+12)\\
 &=\frac{8c^2(c+2)^2}{(c+4)^2}>0.
\end{align*}
Thus \eqref{eq:repeat-slope} is strictly negative.  The leading minor of
$D_\ell$ is negative for every sufficiently large integer $\ell$, so the
Cauchy dual is not subnormal in the double-root case as well.

The unequal-weight case and the three equal-weight cases prove
non-subnormality whenever the two atoms are distinct and non-antipodal.
Combining this with the antipodal theorem of \cite{CGR} proves
Theorem~\ref{thm:main}. \qed

\section{Verification and novelty boundary}

The script \texttt{code/verify\_identities.py} checks, with exact SymPy
arithmetic, the resultant \eqref{eq:resultant}, the exceptional-factor
arguments, the equal-weight factorizations \eqref{eq:Ecommon}, and the
double-root slope and positivity calculation.  The companion numerical
script reconstructs the Fej\'er factor and numerator Gram matrix directly
from the reproducing-kernel algorithm.  It was run on a grid of 29,241
positive-weight/non-antipodal parameter triples; no zero cross Gram value
was found.  Those computations are checks, not substitutes for the proof.

A bounded novelty search through 13 August 2026 used the exact source arXiv
id and title; the phrases ``weighted two atom'', ``arbitrary positive
weights'', and ``Cauchy dual subnormality''; the authors' names; arXiv and
the run's local title/abstract/source indexes; and the citations around
\cite{KPS,CGR}.  No paper settling arbitrary positive weights at two
non-antipodal points was found.  The current arXiv version of \cite{KPS}
still states the extension in its Epilogue.  This is a bounded search, not
an exhaustive priority determination.

The packet settles only measures with exactly two distinct positive atoms.
The classification for three or more support points remains open.  The
recommended review focus is the unequal-weight elimination
\eqref{eq:Rchi}--\eqref{eq:resultant} and the passage from
\eqref{eq:Bn} to \eqref{eq:repeat-slope}.

\begin{thebibliography}{9}
\bibitem{KPS}
M. Khasnis, G. Phatak, and V. Sholapurkar,
\emph{A solution to the Cauchy dual subnormality problem for a cyclic
analytic $2$-isometry with defect operator of rank two},
arXiv:2510.14004v2 (2026).

\bibitem{CGR}
S. Chavan, S. Ghara, and M. R. Reza,
\emph{The Cauchy dual subnormality problem via de Branges--Rovnyak spaces},
arXiv:2103.10059v2 (2021).
\end{thebibliography}

\end{document}
